% ============================================================
%  payrollsMomentum — Technical Function Documentation
%  Quantitative Strategy Designer / SignalBuilder
% ============================================================
\documentclass[11pt, a4paper]{article}

\usepackage[a4paper, margin=2.5cm]{geometry}
\usepackage{amsmath, amssymb, amsthm}
\usepackage{booktabs}
\usepackage{array}
\usepackage{xcolor}
\usepackage{listings}
\usepackage{hyperref}
\usepackage{microtype}
\usepackage{parskip}
\usepackage{graphicx}
\usepackage{float}
\usepackage{enumitem}
\usepackage{caption}
\usepackage{fancyhdr}
\usepackage{titlesec}
\usepackage{mdframed}

% ---------- colours ----------
\definecolor{matlabblue}{RGB}{0,70,127}
\definecolor{codebg}{RGB}{248,248,248}
\definecolor{rulecol}{RGB}{180,180,180}
\definecolor{noteblue}{RGB}{230,240,255}
\definecolor{noteborder}{RGB}{100,140,200}

% ---------- hyperref ----------
\hypersetup{
  colorlinks = true,
  linkcolor  = matlabblue,
  urlcolor   = matlabblue,
  citecolor  = matlabblue,
  pdftitle   = {payrollsMomentum — Technical Documentation},
  pdfauthor  = {SignalBuilder / Quantitative Strategy Designer}
}

% ---------- section formatting ----------
\titleformat{\section}{\large\bfseries\color{matlabblue}}{{\thesection}}{1em}{}[\vspace{2pt}\hrule\vspace{4pt}]
\titleformat{\subsection}{\normalsize\bfseries}{{\thesubsection}}{1em}{}

% ---------- header / footer ----------
\pagestyle{fancy}
\fancyhf{}
\lhead{\small\textcolor{matlabblue}{\texttt{payrollsMomentum} — Technical Documentation}}
\rhead{\small SignalBuilder}
\cfoot{\small\thepage}
\renewcommand{\headrulewidth}{0.4pt}

% ---------- MATLAB listings style ----------
\lstdefinestyle{matlab}{
  language        = Matlab,
  backgroundcolor = \color{codebg},
  basicstyle      = \ttfamily\footnotesize,
  keywordstyle    = \color{matlabblue}\bfseries,
  stringstyle     = \color{orange!70!black},
  commentstyle    = \color{green!50!black}\itshape,
  frame           = single,
  framerule       = 0.4pt,
  rulecolor       = \color{rulecol},
  breaklines      = true,
  columns         = flexible,
  keepspaces      = true,
  showstringspaces= false,
  numbers         = left,
  numberstyle     = \tiny\color{gray},
  numbersep       = 6pt,
}

% ---------- note box ----------
\newmdenv[
  backgroundcolor = noteblue,
  linecolor       = noteborder,
  linewidth       = 1pt,
  innerleftmargin = 8pt,
  innerrightmargin= 8pt,
  innertopmargin  = 6pt,
  innerbottommargin=6pt,
  skipabove       = 6pt,
  skipbelow       = 6pt,
]{notebox}

% ---------- math operators ----------
\DeclareMathOperator{\EMA}{EMA}
\DeclareMathOperator{\Med}{median}
\DeclareMathOperator{\std}{std}
\DeclareMathOperator{\sgn}{sgn}
\newcommand{\dP}{\Delta P}

% ============================================================
\begin{document}

\begin{titlepage}
  \centering
  \vspace*{3cm}
  {\LARGE\bfseries\color{matlabblue} \texttt{payrollsMomentum}}\\[0.6em]
  {\large Technical Function Documentation}\\[2em]
  \hrule height 1pt
  \vspace{1.5em}
  {\large Quantitative Strategy Designer — SignalBuilder}\\[0.4em]
  {\normalsize \texttt{+Model/payrollsMomentum.m}}\\[2em]
  \hrule height 0.4pt
  \vspace{2em}
  {\normalsize
  \begin{tabular}{rl}
    \textbf{Series}       & FRED \texttt{PAYEMS} — Total Nonfarm Payrolls (thousands, SA) \\
    \textbf{Frequency}    & Monthly \\
    \textbf{Output range} & $(-1,\,+1)$ \\
    \textbf{Dependencies} & \texttt{Model.DataManager}, \texttt{Model.BaseModel} \\
    \textbf{MATLAB}       & R2022b+ (built-in \texttt{sqlite}) \\
  \end{tabular}
  }
  \vfill
  {\small\color{gray} Generated: \today}
\end{titlepage}

\tableofcontents
\newpage

% ============================================================
\section{Overview and Motivation}
\label{sec:overview}

\texttt{payrollsMomentum} derives a smooth, bounded momentum signal
from the U.S.\ Bureau of Labor Statistics' Total Nonfarm Payrolls
series (FRED identifier \texttt{PAYEMS}), which is published monthly
and seasonally adjusted. The signal is designed to:

\begin{enumerate}[leftmargin=2em]
  \item Sit strictly inside $(-1,+1)$ at all times, with no discontinuities,
        so it can be consumed directly as a weight or risk-on/off multiplier.
  \item Read strongly positive when payrolls are growing rapidly relative
        to their recent historical norm.
  \item Read near zero (and sign correctly) when growth is decelerating
        or flat.
  \item Read clearly negative when payrolls are declining.
  \item Be robust to three distinct sources of distortion that afflict
        na\"{\i}ve implementations: isolated one-month anomalies
        (weather, strikes, seasonal revisions), regime changes in
        volatility (e.g.\ the post-COVID recovery), and catastrophic
        single-month shocks (e.g.\ April 2020: $-20{,}700$K, roughly
        30 times the magnitude of a typical recession month).
\end{enumerate}

The labour market is a lagging macroeconomic indicator, so the signal
is primarily useful as a \emph{regime filter} — confirming or
challenging risk-on positioning established by leading indicators —
rather than as a high-frequency timing tool.

% ============================================================
\section{Data Source}
\label{sec:data}

\subsection{FRED PAYEMS}

\textbf{Series:} Total Nonfarm Payrolls (All Employees, Thousands of
Persons, Seasonally Adjusted, Monthly). Released by the Bureau of
Labor Statistics on the first Friday of each month for the prior month.

\begin{itemize}[leftmargin=2em]
  \item \textbf{Coverage}: January 1939 -- present ($\sim$1{,}000 observations).
  \item \textbf{Units}: Thousands of jobs.
  \item \textbf{Revisions}: The initial print is revised twice (preliminary
        and final), typically within two months. \texttt{FRED} always
        provides the most recently revised vintage.
  \item \textbf{Typical range}: Level $\approx 60{,}000$K (1970) to
        $158{,}000$K (2024). Monthly change $\dP \in [-21{,}000, +4{,}000]$K,
        with the April 2020 COVID shock being the extreme outlier.
\end{itemize}

\subsection{Fetching and Caching}

Data is retrieved via \texttt{Model.DataManager.fetch\_fred(\textquotesingle PAYEMS\textquotesingle)},
which calls the FRED REST API on the first run and subsequently serves
data from a local SQLite-backed \texttt{.mat} cache. All subsequent
calls within a session — and across sessions until \texttt{update\_data}
is called — are served from cache at zero network cost.

% ============================================================
\section{Signal Construction Pipeline}
\label{sec:pipeline}

Let $P_t$ denote the raw PAYEMS level (thousands of jobs) at month $t$.
The pipeline transforms $P_t$ into a signal $S_t \in (-1, +1)$
through seven sequential steps. Two independent processing paths are
maintained — a \emph{numerator path} and a \emph{denominator path} —
which are combined in the final steps.

\begin{figure}[H]
\centering
\renewcommand{\arraystretch}{1.6}
\begin{tabular}{clcl}
\toprule
\textbf{Step} & \textbf{Operation} & \textbf{Path} & \textbf{Output} \\
\midrule
1 & First difference          & Both   & $\dP_t$ \\
2 & Trailing median prefilter & Num.   & $\tilde{\dP}_t$ \\
3 & Exponential moving average& Num.   & $M_t$ \\
4 & Causal winsorization      & Denom. & $\dP^w_t$ \\
5 & Rolling standard deviation& Denom. & $\sigma_t$ \\
6 & Z-score                   & Merge  & $z_t$ \\
7 & Hyperbolic tangent         & Output & $S_t$ \\
\bottomrule
\end{tabular}
\caption{Signal construction pipeline.}
\end{figure}

% ============================================================
\subsection{Step 1 — First Difference}
\label{sec:step1}

\begin{equation}
  \dP_t \;=\; P_t \;-\; P_{t-1}, \qquad t = 2, 3, \ldots, n
\end{equation}

$\dP_t$ is the net monthly job additions (thousands). It is the
fundamental momentum quantity: positive values indicate job creation,
negative values indicate destruction.

\textbf{Rationale.} Using the level $P_t$ directly would conflate
the secular upward trend in total employment (driven by population
growth) with cyclical acceleration. First-differencing removes the
stochastic trend and isolates the flow of new jobs — the signal
relevant to business cycle positioning.

\subsection{Step 2 — Trailing Median Prefilter (Numerator Path)}
\label{sec:step2}

\begin{equation}
  \tilde{\dP}_t \;=\; \Med\!\bigl(\dP_{t-W_{\mathrm{pre}}+1},\;\ldots,\;\dP_t\bigr),
  \qquad W_{\mathrm{pre}} \geq 1
\end{equation}

where $W_{\mathrm{pre}}$ is controlled by \texttt{PrefilterWindow}
(default 3). The median is computed over trailing (past) observations
only, making it strictly causal.

\textbf{Rationale.} Monthly payroll data contains transient
measurement noise that cannot be explained by genuine labour-market
dynamics:

\begin{itemize}[leftmargin=2em]
  \item \emph{Weather distortions}: Unusually severe winter storms
        suppress construction and leisure hiring for a single month.
  \item \emph{Strike effects}: A major labour action (e.g.\ the UAW
        strike of 2023) temporarily removes workers from the count.
  \item \emph{Seasonal revision anomalies}: BLS benchmark revisions
        create apparent one-month spikes in historical data.
\end{itemize}

The median gate handles these cases with a key property: a single
anomalous print in a window of three is always replaced by the median
of its two neighbours. Two or more months moving in the same
direction pass through unchanged. Contrast this with increasing the
EMA half-life (Step 3): a longer EMA applies proportional additional
lag to \emph{all} movements — genuine trend shifts as well as noise.
The median adds at most one month of lag to genuine regime changes
while completely suppressing isolated prints.

Formally, for a window $\{x_1, x_2, x_3\}$ where $x_1, x_3$ are
``normal'' and $x_2$ is an anomalous spike:
\begin{equation}
  \Med(x_1, x_2, x_3) = x_{\text{middle value}} \approx x_1 \text{ or } x_3
  \quad \text{(not } x_2)
\end{equation}

When \texttt{PrefilterWindow} $= 1$, Step~2 is a no-op:
$\tilde{\dP}_t = \dP_t$.

\subsection{Step 3 — Exponential Moving Average (Numerator Path)}
\label{sec:step3}

\begin{align}
  \alpha &\;=\; 1 - \exp\!\left(-\frac{\ln 2}{\tau}\right) \label{eq:alpha}\\[6pt]
  M_t &\;=\; \alpha\,\tilde{\dP}_t \;+\; (1-\alpha)\,M_{t-1} \label{eq:ema}
\end{align}

where $\tau$ is \texttt{SmoothHalfLife} (default 3 months).

\textbf{Half-life semantics.} The smoothing factor $\alpha$ in
equation~\eqref{eq:alpha} is derived from the half-life condition:
the weight assigned to an observation $\tau$ periods in the past
should equal $\tfrac{1}{2}$ the weight of the current observation.
Expanding the EMA recursion:
\begin{equation}
  M_t \;=\; \sum_{k=0}^{\infty} \alpha\,(1-\alpha)^k\,\tilde{\dP}_{t-k}
\end{equation}
The weight on $\tilde{\dP}_{t-\tau}$ is $\alpha(1-\alpha)^\tau$.
Setting this equal to half the weight on $\tilde{\dP}_t$ (which is
$\alpha$):
\begin{equation}
  (1-\alpha)^\tau = \tfrac{1}{2}
  \;\Longrightarrow\;
  \tau = \frac{-\ln 2}{\ln(1-\alpha)}
  \;\Longrightarrow\;
  \alpha = 1 - e^{-\ln2/\tau}
\end{equation}

This parameterisation is preferable to specifying $\alpha$ directly
because the half-life has an intuitive macroeconomic interpretation:
with $\tau = 3$, a payroll print from six months ago carries only
25\% of the weight of the current print.

\textbf{NaN handling.} If FRED returns a gap in the series (revision
hold-backs, early history), the EMA carries the previous value
forward ($M_t = M_{t-1}$) rather than restarting, preserving
continuity.

\textbf{Rationale.} The EMA acts as a second-stage smoother on top
of the median prefilter. Together, the two-stage filter (median then
EMA) provides stronger noise suppression than either alone with less
aggregate lag than increasing either parameter independently.

\subsection{Step 4 — Causal Expanding-Window Winsorization (Denominator Path)}
\label{sec:step4}

Define the expanding percentile cap at time $t$:
\begin{equation}
  Q_t \;=\; \text{prctile}\!\left(\bigl\{|\dP_s| : s \leq t,\; \dP_s \neq \text{NaN}\bigr\},\;p_w\right)
\end{equation}
where $p_w$ is \texttt{WinsorPct} (default 95). The winsorized series is:
\begin{equation}
  \dP^w_t \;=\; \sgn(\dP_t)\cdot\min\!\bigl(|\dP_t|,\;Q_t\bigr)
\end{equation}

Winsorization is applied only once at least 24 valid observations
have accumulated (to ensure $Q_t$ is estimated from a meaningful sample).

\textbf{Rationale — the COVID problem.} The April 2020 payroll print
($\dP \approx -20{,}700$K) is approximately 30 standard deviations
from the pre-COVID conditional mean. Any rolling standard deviation
computed over a window that contains this observation is massively
inflated. With \texttt{NormWindow} $= 60$ months, \emph{every}
rolling window from 2015 through 2025 includes April 2020, making
the denominator $\sigma_t$ roughly 10--15 times larger than its
pre-COVID value. Post-COVID boom months of $+800$K -- $+1{,}000$K
then produce $z_t \approx 0.1$--$0.2$, collapsing the signal to near
zero despite being genuinely strong job growth.

Winsorizing at the 95th percentile replaces the $-20{,}700$K print
with approximately $-350$K (the 95th percentile of $|\dP|$ computed
over history up to April 2020), keeping $\sigma_t$ in a range
representative of normal-regime volatility.

\textbf{Causality.} $Q_t$ uses only $\{|\dP_s| : s \leq t\}$. No
future data is used, preserving the signal's live-deployment
integrity.

\textbf{Separation of paths.} The winsorization is applied
\emph{only} to the denominator path. The numerator EMA (Step 3)
uses the prefiltered but unwinsorized series. This asymmetry is
deliberate: we want the signal to correctly read \emph{strongly
positive} during a genuine hiring boom — capping the numerator would
suppress the very reading we wish to amplify.

\subsection{Step 5 — Rolling Standard Deviation (Denominator Path)}
\label{sec:step5}

\begin{equation}
  \sigma_t \;=\; \sqrt{\frac{1}{n_t - 1}
    \sum_{s=t-W+1}^{t} \mathbf{1}[\dP^w_s \neq \text{NaN}]
    \bigl(\dP^w_s - \bar{\dP}^w_t\bigr)^2}
\end{equation}

where $W$ is \texttt{NormWindow} (default 60 months), $n_t$ is the
count of non-missing observations in the window, and
$\bar{\dP}^w_t$ is the window mean. This is the sample (Bessel-corrected)
standard deviation, computed by \texttt{movstd} with a strictly
trailing window.

\textbf{Rationale — adaptive normalisation.} Dividing the momentum
estimate $M_t$ by $\sigma_t$ converts it to a dimensionless
\emph{relative} quantity: how many standard deviations of typical
payroll-change volatility does the current smoothed momentum
represent? This is essential for two reasons:

\begin{enumerate}[leftmargin=2em]
  \item \emph{Epoch comparability}: Payroll levels and absolute
        monthly changes have grown with the labour force. A change
        of $+200$K in 1965 (workforce $\approx 60$M) is proportionally
        larger than the same change in 2024 (workforce $\approx 160$M).
        Normalising by rolling $\sigma$ accounts for this scaling.
  \item \emph{Regime sensitivity}: Volatility of monthly changes
        varies significantly across cycles. A 60-month window captures
        one full business cycle, providing a stable but responsive
        reference distribution.
\end{enumerate}

A minimum of $\max(12, \lfloor W/2 \rfloor)$ valid observations is
required; $\sigma_t$ is set to \texttt{NaN} (and the signal is
\texttt{NaN}) during the warm-up period.

\subsection{Step 6 — Z-Score}
\label{sec:step6}

\begin{equation}
  z_t \;=\; \frac{M_t}{\sigma_t}
\end{equation}

$z_t$ is dimensionless and inherits the sign of $M_t$: positive
when payrolls are growing above their recent volatility norm,
negative when declining.

\textbf{Cross-path interaction.} Because $M_t$ (numerator) is
computed from the prefiltered series $\tilde{\dP}$ while $\sigma_t$
(denominator) is computed from the winsorized raw series $\dP^w$,
the two paths serve different roles:

\begin{center}
\renewcommand{\arraystretch}{1.3}
\begin{tabular}{lll}
\toprule
\textbf{Path} & \textbf{Input} & \textbf{Purpose} \\
\midrule
Numerator   & $\tilde{\dP}$ (median prefiltered) & Responsive, noise-gated trend estimate \\
Denominator & $\dP^w$ (winsorized raw)           & Stable, shock-robust volatility reference \\
\bottomrule
\end{tabular}
\end{center}

\subsection{Step 7 — Hyperbolic Tangent Compression}
\label{sec:step7}

\begin{equation}
  \boxed{S_t \;=\; \tanh\!\bigl(\lambda\,z_t\bigr)}
  \qquad S_t \in (-1,\,+1)
\end{equation}

where $\lambda$ is \texttt{TanhScale} (default 1.0).

\textbf{Why \texttt{tanh}?}

$\tanh$ is the natural choice for mapping an unbounded real-valued
signal to a bounded interval. Its key properties for this application are:

\begin{enumerate}[leftmargin=2em]
  \item \textbf{Boundedness}: $\tanh : \mathbb{R} \to (-1,+1)$
        strictly, with no clipping discontinuities.
  \item \textbf{Odd symmetry}: $\tanh(-x) = -\tanh(x)$, so
        positive and negative momentum are treated symmetrically.
  \item \textbf{Linearity near zero}: For small $|z|$,
        $\tanh(\lambda z) \approx \lambda z$, giving fine-grained
        discrimination between modest positive growth and near-flat conditions.
  \item \textbf{Soft saturation}: $\tanh(\lambda z) \to \pm 1$ as
        $z \to \pm\infty$, but never reaches $\pm 1$ for any finite
        $z$. One extreme month cannot push the signal to its maximum;
        sustained directionality is required.
  \item \textbf{Smoothness}: $\tanh$ is infinitely differentiable,
        producing no discontinuities that could cause numerical issues
        in downstream optimisers.
\end{enumerate}

\textbf{Derivative:}
\begin{equation}
  \frac{d}{dz}\tanh(\lambda z) \;=\; \lambda\,\bigl(1 - \tanh^2(\lambda z)\bigr)
\end{equation}
The slope is maximised at $z = 0$ (value $\lambda$) and decays to
zero as $|z| \to \infty$ — exactly the desired behaviour for a
momentum signal that should be \emph{sensitive to changes} when near
neutral and \emph{stable} when already extreme.

\textbf{Effect of \texttt{TanhScale} $\lambda$:}
\begin{equation}
  \lambda \uparrow \;\Rightarrow\; \text{signal saturates faster (more decisive)}
  \qquad
  \lambda \downarrow \;\Rightarrow\; \text{signal stays more graded near extremes}
\end{equation}

At $\lambda = 1$, $z = 1$ maps to $\tanh(1) \approx 0.76$. This
means a smoothed payroll print one historical standard deviation
above the rolling mean produces a signal of about $+0.76$ — clearly
positive but not at the ceiling, leaving room to strengthen further.

% ============================================================
\section{Calibration}
\label{sec:calibration}

Table~\ref{tab:calib} gives approximate signal values for
representative payroll regimes under default parameters
(\texttt{PrefilterWindow}$=3$, \texttt{SmoothHalfLife}$=3$,
\texttt{NormWindow}$=60$, \texttt{TanhScale}$=1.0$, \texttt{WinsorPct}$=95$).
The rolling standard deviation of winsorized $\dP$ over the full
post-1970 sample is approximately $\sigma \approx 130$K.

\begin{table}[H]
\centering
\renewcommand{\arraystretch}{1.35}
\begin{tabular}{lrrrrl}
\toprule
\textbf{Regime} & $\dP_t$ (K) & $M_t$ (K) & $\sigma_t$ (K) & $z_t$ & $S_t \approx$ \\
\midrule
Post-COVID hiring boom     & $+900$  & $+750$  & 130 & $+5.8$ & $+0.99$ \\
Strong expansion           & $+400$  & $+360$  & 130 & $+2.8$ & $+0.99$ \\
Solid expansion            & $+200$  & $+185$  & 130 & $+1.4$ & $+0.89$ \\
Moderate growth            & $+100$  & $+95$   & 130 & $+0.7$ & $+0.61$ \\
Stalling (near zero)       & $+20$   & $+20$   & 130 & $+0.2$ & $+0.16$ \\
Flat / no-growth           & $0$     & $0$     & 130 & $0.0$  & $0.00$  \\
Mild contraction           & $-100$  & $-95$   & 130 & $-0.7$ & $-0.61$ \\
Recession                  & $-300$  & $-270$  & 130 & $-2.1$ & $-0.97$ \\
Isolated 1-month spike     & $\pm500$ & $\approx 0$ & 130 & $\approx 0$ & $\approx 0.00$ \\
\bottomrule
\end{tabular}
\caption{Approximate signal values under default parameters.
$M_t$ shown as EMA-converged value after several months at the stated $\dP$.
The last row illustrates the median prefilter's effect on a single-month anomaly.}
\label{tab:calib}
\end{table}

% ============================================================
\section{Parameter Reference}
\label{sec:params}

\begin{table}[H]
\centering
\renewcommand{\arraystretch}{1.5}
\begin{tabular}{p{3.8cm} p{1.4cm} p{1.4cm} p{7.4cm}}
\toprule
\textbf{Parameter} & \textbf{Type} & \textbf{Default} & \textbf{Description and guidance} \\
\midrule
\texttt{PrefilterWindow} &
  integer $\geq 1$ &
  3 &
  Length $W_{\mathrm{pre}}$ of the trailing median window applied to $\dP$ before the EMA. \newline
  1 = disabled; 3 = removes isolated single-month anomalies; 5 = removes up to two consecutive anomalous months. Larger values reduce noise but increase the lag for genuine inflection points by at most $\lfloor W_{\mathrm{pre}}/2 \rfloor$ months. \\

\texttt{SmoothHalfLife} &
  double $> 0$ &
  3 &
  EMA half-life $\tau$ in months. Controls how quickly the smoothed momentum $M_t$ adapts to new prints.\newline
  Lower ($\tau=2$): faster, noisier. Higher ($\tau=6$): slower, smoother. Increasing this is \emph{not} a substitute for \texttt{PrefilterWindow}: it adds lag to genuine signals proportionally. \\

\texttt{NormWindow} &
  integer $\geq 12$ &
  60 &
  Trailing window width $W$ (months) for computing $\sigma_t$ from winsorized $\dP^w$.\newline
  60 months $\approx$ one business cycle: stable enough to smooth regime transitions, short enough to respond to structural shifts in payroll volatility over a decade. \\

\texttt{TanhScale} &
  double $> 0$ &
  1.0 &
  Scale parameter $\lambda$ applied to $z_t$ before $\tanh$. Stretches or compresses the signal's dynamic range.\newline
  $\lambda > 1$: signal saturates closer to $\pm 1$ for moderate $z$ — more binary/decisive.\newline
  $\lambda < 1$: signal stays graded across a wider range of $z$ — more analogue. \\

\texttt{WinsorPct} &
  double $\in [0, 100)$ &
  95 &
  Expanding-window percentile $p_w$ used to compute the winsorization cap $Q_t$ for the denominator path.\newline
  0 = no winsorization (original behaviour). 95 caps the top 5\% of $|\dP|$ observations cumulatively seen up to $t$. Increase to 97--98 for a lighter touch; decrease toward 90 for more aggressive shock-proofing of $\sigma_t$. \\

\texttt{StartDate} &
  string &
  \texttt{""} &
  Optional FRED fetch start date in \texttt{'yyyy-MM-dd'} format. Empty string requests the full PAYEMS history (January 1939 -- present). Note: signal warm-up requires at least \texttt{NormWindow}$+2$ valid observations before $S_t$ is non-NaN. \\
\midrule
\texttt{dm} &
  \texttt{DataManager} &
  (required) &
  Initialised \texttt{Model.DataManager} handle. Must have \texttt{init()} called and \texttt{FredApiKey} set. \\
\bottomrule
\end{tabular}
\caption{Parameter reference for \texttt{payrollsMomentum}.}
\label{tab:params}
\end{table}

% ============================================================
\section{Tuning Guide}
\label{sec:tuning}

\subsection{Reducing Noise Without Sacrificing Reactivity}

\begin{itemize}[leftmargin=2em]
  \item Increase \texttt{PrefilterWindow} from 3 to 5. This suppresses
        up to two consecutive anomalous prints while adding at most
        2 months of lag to genuine inflections.
  \item \emph{Avoid} compensating for noise by increasing
        \texttt{SmoothHalfLife} beyond 4--5: the resulting lag on
        genuine recessions and recoveries becomes macroeconomically
        significant.
\end{itemize}

\subsection{Making the Signal More Decisive}

\begin{itemize}[leftmargin=2em]
  \item Increase \texttt{TanhScale} to 1.3--1.5. Solid expansion
        months ($+150$K, $z \approx 1.1$) then map to
        $\tanh(1.5 \times 1.1) \approx 0.94$ rather than $0.80$.
  \item The signal still cannot reach $\pm 1$ for any finite input.
\end{itemize}

\subsection{Post-Shock Recovery}

\begin{itemize}[leftmargin=2em]
  \item If the signal is visibly muted for several years following a
        shock event, increase \texttt{WinsorPct} from 95 to 97.
        This applies a lighter cap, preserving more of the historical
        volatility in $\sigma_t$ while still preventing the catastrophic
        outlier from dominating.
  \item Do not increase \texttt{WinsorPct} beyond 99: at that point
        the protection against genuine extreme shocks is lost.
\end{itemize}

\subsection{Faster Reaction to Trend Shifts}

\begin{itemize}[leftmargin=2em]
  \item Reduce \texttt{SmoothHalfLife} from 3 to 2. With
        $\alpha = 1 - e^{-\ln2/2} \approx 0.293$, a genuine
        multi-month payroll shift reaches 75\% of its full signal
        strength roughly one month sooner.
  \item Consider also reducing \texttt{PrefilterWindow} from 3 to 1
        (disabled) if the additional half-month of lag introduced by
        the median gate is unacceptable in your application.
\end{itemize}

% ============================================================
\section{Outputs}
\label{sec:outputs}

The function returns a \texttt{timetable} with \texttt{RowTimes}
preserved from the FRED series (UTC timezone, first day of month):

\begin{center}
\renewcommand{\arraystretch}{1.3}
\begin{tabular}{llll}
\toprule
\textbf{Variable} & \textbf{Units} & \textbf{Range} & \textbf{Description} \\
\midrule
\texttt{Payrolls}   & K jobs     & $[60{,}000,\,160{,}000]$ & Raw PAYEMS level \\
\texttt{dPayrolls}  & K jobs/mo  & $(-\infty,\,+\infty)$     & Raw MoM change \\
\texttt{SmoothedMom}& K jobs/mo  & $(-\infty,\,+\infty)$     & Prefiltered + EMA momentum ($M_t$) \\
\texttt{Signal}     & —          & $(-1,\,+1)$               & Final bounded signal ($S_t$) \\
\bottomrule
\end{tabular}
\end{center}

The timetable \texttt{Description} property encodes the parameters
used to generate it:
\begin{lstlisting}[style=matlab, numbers=none]
>> tt.Properties.Description
ans = 'PAYEMS momentum | PrefilterWin=3, SmoothHL=3 mo, NormWin=60 mo, TanhScale=1, WinsorPct=95'
\end{lstlisting}

% ============================================================
\section{Dependencies and Portability}
\label{sec:deps}

\subsection{Required Files}
To run \texttt{payrollsMomentum} outside the SignalBuilder project,
copy only:
\begin{itemize}[leftmargin=2em]
  \item \texttt{+Model/payrollsMomentum.m}
  \item \texttt{+Model/DataManager.m}
  \item \texttt{+Model/BaseModel.m}
\end{itemize}

\subsection{MATLAB Requirements}
\begin{itemize}[leftmargin=2em]
  \item MATLAB R2022b or later (built-in \texttt{sqlite} function).
        R2021b requires the Database Toolbox for \texttt{sqlite}.
  \item No Statistics Toolbox: \texttt{prctile} has been replaced by
        the local helper \texttt{payrollsMomentum\_prctile} (nearest-rank
        method, implemented with \texttt{sort} and index arithmetic).
  \item Network access on the first run for the FRED API call.
        Subsequent runs are served from local cache at zero network cost.
\end{itemize}

% ============================================================
\section{Usage Examples}
\label{sec:usage}

\begin{lstlisting}[style=matlab, caption={Minimal usage}]
dm = Model.DataManager();
dm.init();
tt = Model.payrollsMomentum(dm);
\end{lstlisting}

\begin{lstlisting}[style=matlab, caption={Custom parameters}]
tt = Model.payrollsMomentum(dm, ...
    PrefilterWindow = 5,   ... % stronger noise gate
    SmoothHalfLife  = 2,   ... % faster reaction
    TanhScale       = 1.3, ... % more decisive signal
    WinsorPct       = 97,  ... % lighter shock cap
    StartDate       = "1990-01-01");
\end{lstlisting}

\begin{lstlisting}[style=matlab, caption={Compare prefiltered vs.\ no prefilter}]
tt_raw  = Model.payrollsMomentum(dm, PrefilterWindow=1);
tt_filt = Model.payrollsMomentum(dm, PrefilterWindow=3);

figure;
yyaxis left
bar(tt_filt.Time, tt_filt.dPayrolls, 'FaceAlpha', 0.3)
ylabel('MoM change (K jobs)')

yyaxis right
plot(tt_raw.Time,  tt_raw.Signal,  ':', 'LineWidth', 1, ...
    'DisplayName', 'No prefilter')
hold on
plot(tt_filt.Time, tt_filt.Signal, '-', 'LineWidth', 2, ...
    'DisplayName', 'PrefilterWindow=3')
yline(0, 'k--'); ylim([-1.1 1.1])
ylabel('Signal'); legend; grid on
title(tt_filt.Properties.Description)
\end{lstlisting}

\begin{lstlisting}[style=matlab, caption={Incremental cache refresh}]
dm.update_data('PAYEMS');          % only fetches the delta from last cache date
tt = Model.payrollsMomentum(dm);   % immediately uses refreshed cache
\end{lstlisting}

% ============================================================
\section{Mathematical Summary}
\label{sec:summary}

For reference, the complete signal in closed form (suppressing NaN
handling and warm-up logic):

\begin{align}
  \dP_t &= P_t - P_{t-1} \\[4pt]
  \tilde{\dP}_t &= \Med_{s \in [t-W_{\mathrm{pre}}+1,\,t]}\!\bigl(\dP_s\bigr) \\[4pt]
  M_t &= \alpha\,\tilde{\dP}_t + (1-\alpha)\,M_{t-1},
    \quad \alpha = 1 - e^{-\ln2/\tau} \\[4pt]
  Q_t &= \text{prctile}\!\left(\bigl\{|\dP_s|\bigr\}_{s \leq t},\; p_w\right) \\[4pt]
  \dP^w_t &= \sgn(\dP_t)\cdot\min\!\bigl(|\dP_t|,\;Q_t\bigr) \\[4pt]
  \sigma_t &= \std_{s \in [t-W+1,\,t]}\!\bigl(\dP^w_s\bigr) \\[4pt]
  z_t &= \frac{M_t}{\sigma_t} \\[6pt]
  S_t &= \tanh\!\bigl(\lambda\,z_t\bigr) \;\in\; (-1,\,+1)
\end{align}

\textbf{Parameters:}
$W_{\mathrm{pre}}$ (\texttt{PrefilterWindow}),
$\tau$ (\texttt{SmoothHalfLife}),
$W$ (\texttt{NormWindow}),
$\lambda$ (\texttt{TanhScale}),
$p_w$ (\texttt{WinsorPct}).

\vfill
\hrule
\vspace{6pt}
{\small\color{gray}
\textbf{File}: \texttt{+Model/payrollsMomentum.m} \quad
\textbf{Documentation}: \texttt{+Model/payrollsMomentum\_doc.pdf} \quad
\textbf{Project}: SignalBuilder / Quantitative Strategy Designer
}

\end{document}
